\section{Background}
\label{sec:background}

This section situates the Recursive Spawning Protocol (RSP) within the broader landscape of prior work across five domains: agent lifecycle management in multi-agent systems, accountability and provenance in distributed AI agent ecosystems, fixed versus mutable delegation chains, hierarchical task decomposition in AI, and the BX3 Framework as the architectural substrate that makes immutable parent chains enforceable at runtime.

% ------------------------------------------------------------------
\subsection{Agent Lifecycle Management in Multi-Agent Systems}
\label{sec:lifecycle}

Multi-agent systems (MAS) composed of LLM-powered agents introduce lifecycle challenges fundamentally different from conventional distributed systems. Each agent is a stateful reasoning entity capable of planning, tool use, adaptive behavior, and recursive self-modification, making lifecycle management substantially more complex than traditional microservice orchestration~\cite{moore2025hmas}. Agent lifecycle in this context encompasses creation, authorization, execution, communication, suspension, and decommissioning of entities that may multiply in real time as orchestrator agents spawn sub-agents in response to emerging subtask demands.

The defining failure mode in unmanaged agent fleets is \emph{agent sprawl}: the proliferation of agents without clear ownership, explicit lifecycle boundaries, or consistent policy enforcement at handoff boundaries~\cite{moore2025hmas}. AgentOrchestra identifies a complementary concern: agents spawned for specialized subtasks frequently fail to inherit the policy context of their parent, creating enforcement gaps precisely at handoff boundaries~\cite{agentorchestra2025}. PROV-AGENT extends the W3C PROV standard into agentic workflows, capturing agent-centric metadata (prompts, responses, decisions) and linking them to the broader workflow provenance graph, demonstrating that end-to-end provenance enables post-hoc fault isolation in heterogeneous multi-agent pipelines~\cite{prov-agent2025}. The Agent Lifecycle Protocol (ALP) working group formalizes lifecycle stages---provisioning, activation, monitoring, and retirement---alongside drift detection mechanisms that flag deviations from approved behavior~\cite{alp2025}. Together these works establish that lifecycle management is not separable from accountability: a system cannot reconstruct \emph{which agent performed which action, under whose authorization, at what time, with what effect} without explicit lifecycle infrastructure baked in at the architectural level.

% ------------------------------------------------------------------
\subsection{Accountability and Provenance in Distributed Agent Ecosystems}
\label{sec:accountability}

The accountability deficit in distributed AI systems differs in kind from conventional software accountability. When an LLM-powered agent makes a decision inside a multi-agent chain, the decision may emerge from complex, opaque reasoning over contextual state that resists post-hoc reconstruction~\cite{dzone2025accountability}. Without explicit provenance infrastructure, it is functionally impossible to attribute outcomes to responsible actors, verify policy compliance at each delegation hop, or detect scope drift when a sub-agent acts beyond its authorized mandate.

SentinelAgent formalizes this problem through the Delegation Chain Calculus (DCC), which defines seven properties of sound delegation chains: six deterministic (authority narrowing, policy preservation, forensic reconstructibility, cascade containment, scope-action conformance, output schema conformance) and one probabilistic (intent preservation)~\cite{sentinelagent}. DCC's \emph{forensic reconstructibility} property demands that every delegation hop be recorded immutably, not merely accessible via a mutable log. SentinelAgent implements these properties at runtime via a non-LLM Delegation Authority Service, achieving a 100\% true-positive rate at 0\% false-positive on DelegationBench across 516 scenarios; properties P3--P7 are mechanically verified using TLA+ across 2.7 million states with zero violations~\cite{sentinelagent}. The key lesson for RSP is that forensic reconstructibility is satisfiable only if delegation records are append-only by construction.

The Human Delegation Provenance (HDP) Protocol provides a cryptographic implementation of this principle: a token binds a human principal to their declared authorization scope and session identifier, then each agent appends a cryptographically signed hop to an append-only chain, with offline verification requiring only the issuer's Ed25519 public key and the current session identifier---no central registry or third-party trust anchors~\cite{hdp-draft}. Chain of Consciousness extends this to an append-only SHA-256 hash chain of lifecycle events anchored to Bitcoin via OpenTimestamps, with W3C DID-based persistent agent identity enabling continuity proofs that link discontinuous sessions into a single verifiable identity arc~\cite{chainofconsciousness2026}. Layered Mutability provides the conceptual model for distinguishing mutable from immutable elements in persistent agent systems, showing that selectively fixing certain structural elements while permitting controlled mutation elsewhere achieves a favorable tradeoff between adaptability and auditability~\cite{layeredmutability2026}. Warrant Certificate Authorities (WCA) propose a PKI-style trust hierarchy for AI-agent tool-call attestations, with graduated attestation levels analogous to SLSA build levels, providing tamper-evident certification of data provenance across agent tool chains~\cite{wca-draft}.

% ------------------------------------------------------------------
\subsection{Fixed versus Mutable Delegation Chains}
\label{sec:delegation-chains}

The tension between fixed and mutable delegation chains is a central unresolved question in agentic system design. A \emph{fixed delegation chain} establishes delegation relationships at spawn time and preserves them immutably throughout execution: the parent pointer cannot be severed or redirected without detection, and any attempt to do so constitutes a verifiable policy violation. A \emph{mutable delegation chain} permits dynamic re-delegation, agent migration, or role reassignment during execution, offering flexibility at the cost of forensic traceability.

Mutable delegation is the dominant approach in current LLM multi-agent frameworks. AgentSpawn adaptively recruits collaborators based on runtime task complexity but spawned agents carry no cryptographic record of their parent's identity beyond a shared conversation context~\cite{agentspawn2025}. FIRMHIVE spawns hierarchical agent trees for firmware security analysis but relies on in-memory parent references that are lost on process restart~\cite{firmhive2025}. In both cases, post-hoc accountability requires trusting the system's internal log integrity or reconstructing behavior from partial, mutable records. AgentOrchestra confirms that spawned agents may not inherit policy context from their parent, creating enforcement gaps at every handoff boundary~\cite{agentorchestra2025}. These are not incidental oversights; they are structural consequences of treating delegation topology as a runtime optimization rather than an immutable authorization record.

The accountability cost of mutable chains is severe. When a parent reference can change after the fact, the chain of responsibility becomes history-dependent in a way that defeats audit: an action taken at time $t$ may be attributable to one parent, but retroactive pointer updates rewrite that attribution without detection. SentinelAgent identifies specific adversaries enabled by mutable chains: a Privilege Escalator (Type A4) exploits mid-chain mutability to expand its own authority, and a Chain Poisoner (Type A5) corrupts delegation records to make unauthorized actions appear authorized~\cite[Table I]{sentinelagent}. HDP and SentinelAgent independently conclude that immutable delegation with cryptographic identity binding is the required alternative for accountable systems~\cite{hdp-draft,sentinelagent}. BX3's Immutable Parent Pointer resolves the fixed-versus-mutable tension architecturally: the pointer itself is immutable once established, enforced by the Fact Layer, while the content of what is delegated may be bounded, reinterpreted, and escalated through the delegation chain under controlled conditions defined by the Bounds Engine.

% ------------------------------------------------------------------
\subsection{Hierarchical Task Decomposition in AI}
\label{sec:decomposition}

Hierarchical task decomposition (HTD) addresses the fundamental scalability limit of flat agent architectures: a single agent cannot efficiently represent or execute the full complexity of a high-level objective over extended time horizons. HTD recursively breaks high-level objectives into subgoals handled by specialized agents, with a coordination layer managing execution order, data flow, and termination conditions.

ReAcTree introduces hierarchical agent trees with two node types: control-flow nodes that organize execution order, and agent nodes that reason, act, and expand into further subgoals~\cite{reactree2025}. Dynamic tree expansion allows agent nodes to autonomously generate and refine subgoals at runtime, substantially outperforming flat sequential baselines (61\% versus 31\% goal success on WAH-NL using Qwen 2.5 72B). StructuredAgent combines online hierarchical planning using dynamic AND/OR trees with a structured memory module that tracks candidate solutions and constraints, enabling efficient exploration of contingent action sequences when a path fails~\cite{structuredttt2026}. DeepAgent uses a Hierarchical Task DAG (HTDAG) to dynamically decompose high-level objectives into interdependent sub-tasks across multiple layers, with a recursive two-stage planner-executor enabling continuous task refinement and real-time graph modification including user-directed intervention~\cite{deepagent2025}. FIRMHIVE operationalizes this through a runtime Tree of Agents (ToA) where each node performs self-contained executable primitives, achieving approximately 16$\times$ more reasoning steps and 2.3$\times$ broader file exploration than baseline LLM-agent systems on firmware security analysis tasks~\cite{firmhive2025}.

AGENTHIVE treats delegation as a first-class primitive: any agent can recursively spawn sub-agents, forming dynamic topologies (trees, forests, star structures) driven by task demands rather than pre-defined graphs, improving scalability and adaptability for open-ended tasks compared to static MAS architectures~\cite{agenthive2025}. The Recursive Delegation Engine (RDE) provides a formal decision-theoretic treatment of HTD: agents use recursive decision rules to choose between executing, decomposing, or delegating subtasks, with success probabilities propagated back up the chain to inform parent decisions~\cite{rde2025}. RDE's key insight is that the delegation graph itself is a learned structure, but RDE does not address the accountability implications of adaptive delegation decisions, treating them as optimization problems rather than as authorization events requiring provenance.

These works collectively establish that hierarchical, recursive decomposition is both necessary for long-horizon task competence and empirically validated across diverse domains. The outstanding challenge---which RSP addresses---is integrating this decomposition capability with immutable, accountable delegation chains from the outset rather than adding them as an afterthought.

% ------------------------------------------------------------------
\subsection{The BX3 Framework as Substrate for Immutable Parent Chains}
\label{sec:bx3-substrate}

The BX3 Framework organizes any complex autonomous system into three immutable functional layers: the \emph{Purpose Layer} (intent, judgment, and upstream accountability anchored to a human principal), the \emph{Bounds Engine} (bounded reasoning, constrained execution, and controlled reinterpretation), and the \emph{Fact Layer} (deterministic enforcement, hard constraint, and forensic auditability)~\cite{bx3-framework}. The framework is actor-agnostic: any entity capable of satisfying a layer's functional requirements may occupy that layer---human, AI system, mechanical process, or hybrid composition. What cannot vary are the properties each layer must maintain.

The upstream accountability guarantee is the property most directly relevant to RSP: when any node encounters a state it cannot resolve within its bounds, accountability escalates recursively upward through the system hierarchy, bypassing all machine actors, until it reaches a human anchor. The property \emph{fails upward into human consciousness---never downward into algorithmic chaos} is not a policy directive but an architectural constraint enforced by the Fact Layer~\cite{bx3-framework}. This guarantee previously applied at the system level; RSP extends it to the \emph{per-node} level.

BX3's Immutable Parent Pointer converts the system-level accountability guarantee into a per-node architectural property. Every spawn event is treated as a delegation event with three formally defined components: a parent identity (immutable, enforced by the Fact Layer pre-commit hook), a bounded mandate derived from the parent's Purpose Layer authorization (enforced by the Bounds Engine), and an immutable record in the Forensic Ledger (enforced by the Fact Layer). This converts the BX3 three-layer model from a system-level accountability architecture into a \emph{per-node accountability architecture}, enabling recursive spawning to preserve upstream accountability at every tier of the task decomposition hierarchy.

Critically, the Bounds Engine's bounded reinterpretation privilege ensures that a child node's Bounds Engine may reason within the mandate established by its parent but cannot expand that mandate unilaterally. If the child encounters a state outside its bounded mandate, it escalates to its parent via the same upstream accountability mechanism. This creates a recursive accountability chain in which every node at every depth has a formally accountable parent, terminating at the human anchor at the root. The provenance of every decision is traceable to a human authorization event.

Recent complementary work confirms the timeliness of this contribution. SentinelAgent's DCC independently identifies forensic reconstructibility and scope-action conformance as required properties of accountable delegation~\cite{sentinelagent}. PROV-AGENT demonstrates that W3C PROV semantics can be extended to capture agent-centric metadata in multi-agent workflows~\cite{prov-agent2025}. HDP provides the cryptographic primitives for append-only hop records~\cite{hdp-draft}. Layered Mutability provides the conceptual model for distinguishing mutable from immutable elements~\cite{layeredmutability2026}. What none of these works supplies---and what the BX3 Framework with RSP provides---is a unified, actor-agnostic architectural substrate that simultaneously enforces immutability of parent pointers, bounded reinterpretation by spawned nodes, and recursive escalation of unresolved states to a human anchor, all as first-class, formally specified properties rather than as implementation heuristics.
